% !PRINTED 624
\weberpage{624}

\[
\begin{array}{lcccccc}
 & a_1\sqrt{\Delta}, & a_2, & a_3\sqrt{\Delta}, & a_4, & a_5\sqrt{\Delta}, & a_6\\
\text{the degrees} & 2 & 4 & 6 & 8 & 10 & 12.
\end{array}
\]
But $\sqrt{\Delta}$ is of degree 10 in the $x$, and consequently $a_1=0$,
$a_3=0$ must hold, and $a_5$ must be a number.

One can compute the coefficients $a_2$, $a_4$, $a_5$, $a_6$ by actual
formation of the expressions according to the rules for the representation of the symmetric
functions, which has been done by \gesperrt{Cayley}. But the calculation and the formulae are very
lengthy.

We shall content ourselves here with forming the resolvent for
a special case. The value found in doing so
for the number $a_5$ is then of course generally valid.

Let the equation of degree 5 have the \gesperrt{Bring-Jerrard}
form\footnote{\gesperrt{Runge}, Acta mathematica, Bd.\,7.}
\weq{13}{x^5+\alpha x+\beta=0.}
Then $a_2$, $a_4$, $a_6$ are integral rational functions of $\alpha$ and $\beta$.
Of the coefficients of the general equation (6), $d$ is of
the fourth, $e$ of the fifth order in the $x$. From this it follows that
in $a_2$, $a_4$, $a_6$ the coefficient $e$ can occur only multiplied by one of the coefficients
$a$, $b$, $c$, and hence that, if $a$, $b$, $c$
are set to zero, $e$ must drop out of these expressions.
For equation (13), therefore, $a_2$, $a_4$, $a_6$ cannot depend
on $\beta$, and, having regard to the degree, one obtains, if
$m$, $m_1$, $m_2$, $m_3$ denote numbers,
\weq{14}{a_2=m_1\alpha,\quad a_4=m_2\alpha^2,\quad a_6=m_3\alpha^3,\quad a_5=m.}

The numbers $m$, $m_1$, $m_2$, $m_3$ can be determined from a special
assumption. If we set $\beta=0$, then
\weq{15}{\begin{aligned}
x_0&=0,\quad x_1=\sqrt[4]{-\alpha},\quad x_2=i\sqrt[4]{-\alpha},\\
&\hphantom{{}=0,\quad{}} x_3=-\sqrt[4]{-\alpha},\quad x_4=-i\sqrt[4]{-\alpha},
\end{aligned}}
and from this there results
\weq{16}{\begin{aligned}
\sqrt{\Delta} &= (x_0-x_1)(x_0-x_2)(x_0-x_3)(x_0-x_4)(x_1-x_2)\\
&\quad (x_1-x_3)(x_1-x_4)(x_2-x_3)(x_2-x_4)(x_3-x_4)\\
&= 16i\sqrt{-\alpha}^5 = -16\sqrt{\alpha}^5\\
\Delta &= 256\alpha^5.
\end{aligned}}
